%% 7-11 行列の使いみち（読みもの。第14回：組 V7-6。ページ制限なし。偶数ページ）
%% 読みものの型：現象の場面→どこで行列が働いているか→専門書での書き方。問題は置かない。
%% ほかの巻・ほかの章が指している形：
%%   委員会の決定6：四端子定数（送電線・分布定数回路、6-11 の cosh・sinh）と逆運動学（ヤコビ行列の逆行列で少しずつ近づける）を入れる。
%%   6-back の予告「ばね・電気・熱の相似の表（6-15）」→ ④「ばねの網と抵抗の網は同じ形の行列」（6-15 の表とは組み方がちがうことを書く）。
%%   5-6:666「第7巻」の回収（ロボットの腕の速さ）→ ①②。2リンクの手は 7-6 と同じ (0.325, 0.391)。
%%   7-3 の箱の行列（出口の組を入れると入口の組が出てくる向き。成分は数で書き、電気の本では A〜D）、7-4 の「箱の行列の行列式は1」。
%%   7-10 の (I+AΔt)^n → e^{At}（名前だけ）を、長い電線の「小さな箱を n 個つなぐ」で回収。
\makeatletter
% 本文の小見出し（太字の1行）
\newcommand\kviilhead[1]{\par\addvspace{6pt}\Needspace{6\baselineskip}\noindent{\sffamily\bfseries\color{\bkTT}#1}\par\nobreak\vspace{2pt}}
% 囲みの中の小見出し（それだけで1行）
\newcommand\kviilbox[1]{\par\addvspace{3pt}\noindent\textbf{#1}\par\nobreak\vspace{1pt}\noindent\ignorespaces}
% 大事な式のわく
\newcommand\kviilframe[1]{\par\addvspace{4pt}{\centering\tikz\node[draw=\bkTT, line width=0.9pt, fill=bkTheme!10, rounded corners=1.5mm, inner xsep=6mm, inner ysep=2.5mm, align=center]{#1};\par}\addvspace{4pt}\noindent\ignorespaces}
% 横向きの抵抗（まん中が (#1,#2)）
\newcommand\kviilres[2]{\filldraw[fill=white, draw=bkInk, line width=0.8pt] (#1-4,#2-1.6) rectangle (#1+4,#2+1.6);}
% たて向きの抵抗
\newcommand\kviilresv[2]{\filldraw[fill=white, draw=bkInk, line width=0.8pt] (#1-1.6,#2-4) rectangle (#1+1.6,#2+4);}
% ばね（横向き）
\newcommand\kviilspring[3]{\draw[bkInk, line width=0.8pt, decorate, decoration={zigzag, amplitude=1.3mm, segment length=2mm, pre length=1.5mm, post length=1.5mm}] (#1,#3) -- (#2,#3);}
% ページ（検索の例）  \kviilpage{(x,y)}{名前}{色}
\newcommand\kviilpage[3]{\node[draw=#3!80!black, fill=#3!15, line width=0.9pt, minimum width=14mm, minimum height=17mm, rounded corners=1mm, font=\sffamily\bfseries\large, text=bkInk] at #1 {#2};}
\makeatother

\bkchapter{7}{11}{行列の使いみち}{\bklevelrange{3}{5}}{行列の使いみち}
\bkmichishirube{6-11,6-15,7-3,7-5,7-7,7-8,7-10}{8-3,9-10,9-11,9-12}

\begin{goal}
\begin{itemize}[label={\color{\bkTT}●}, leftmargin=1.6\zw, itemsep=0pt, topsep=1pt]
\item \bklevel{3}\enspace ロボット・電気回路・建物・画像・検索で、どの章の行列が働いているかを言える。
\item \bklevel{4}\enspace ばねの網と抵抗の網、箱をつないだ回路が、同じ形の行列で書けることを説明できる。
\item \bklevel{5}\enspace 逆運動学と長い電線（分布定数回路）を、「行列を何回も使う計算」として読める。
\end{itemize}
\end{goal}

\begin{hatena}
\noindent\begin{minipage}[c]{0.5\linewidth}
\raggedright
工場のロボット、家のコンセントにつながる送電線、スマホの写真、地震に強いビル、インターネットの検索。どれも、たくさんの数を組にして、同じルールでいっせいに計算しています。\par\vspace{4pt}
行列は、どこで、どんな仕事をしているのでしょう。
\end{minipage}\hfill
\begin{minipage}[c]{0.46\linewidth}
\centering
\begin{tikzpicture}[x=1mm, y=1mm]
  % ロボット
  \fill[bkGray!40] (0,0) rectangle (12,2);
  \draw[bkOrange!80!black, line width=2.4pt, line cap=round] (6,2) -- (10,12) -- (19,14);
  \filldraw[fill=white, draw=bkInk, line width=0.7pt] (6,2) circle (1.1) (10,12) circle (1.1);
  \node[font=\sffamily\footnotesize, anchor=north] at (8,-1.5) {ロボット};
  % 送電線の鉄塔
  \begin{scope}[shift={(34,0)}]
    \draw[bkInk, line width=0.8pt] (-4,0) -- (0,18) -- (4,0) (-2.7,6) -- (2.7,6) (-1.6,11) -- (1.6,11) (-6,14) -- (6,14);
    \draw[bkBlue!70!black, line width=0.6pt] (-6,14) .. controls (-12,10) .. (-16,13) (6,14) .. controls (12,10) .. (16,13);
    \node[font=\sffamily\footnotesize, anchor=north] at (0,-1.5) {送電線};
  \end{scope}
  % 写真（画素）
  \begin{scope}[shift={(0,-30)}]
    \foreach \i in {0,...,3} {\foreach \j in {0,...,3} {\pgfmathsetmacro\g{15+15*mod(\i+2*\j,4)}\fill[bkInk!\g] ({2+4*\i},{2+4*\j}) rectangle ++(3.6,3.6);}}
    \node[font=\sffamily\footnotesize, anchor=north] at (9.8,0.5) {写真};
  \end{scope}
  % ビル
  \begin{scope}[shift={(34,-30)}]
    \filldraw[fill=bkBlue!12, draw=bkBlue!70!black, line width=0.7pt] (-6,2) rectangle (6,20);
    \foreach \y in {6,10,14} {\draw[bkBlue!70!black, line width=0.5pt] (-6,\y) -- (6,\y);}
    \foreach \x in {-3,1} {\foreach \y in {3,7,11,15} {\fill[bkBlue!40] (\x,\y) rectangle ++(2,2);}}
    \node[font=\sffamily\footnotesize, anchor=north] at (0,0.5) {ビル};
  \end{scope}
\end{tikzpicture}
\end{minipage}
\end{hatena}

\begin{omoidasou}
\bkrecall{7-2}{ヤコビ行列 $J$：関節の角速度を入れると、手の速度が出てくる機械（$\dot{\vec p}=J\dot{\vec\theta}$）。}
\bkrecall{7-3}{電気の箱の行列は、出口の組 $(V_2,\ I_2)$ を入れると入口の組 $(V_1,\ I_1)$ が出てくる。箱をつなぐと行列のかけ算。}
\bkrecall{7-4}{行列式が0の行列は、平らにつぶす（逆行列がない）。箱の行列の行列式は1。$\det J=L_1L_2\sin\theta_2$。}
\bkrecall{7-5}{$A\vec x=\vec b$ は、$\vec x=A^{-1}\vec b$ で解ける。コンピュータは消去法で解く。}
\bkrecall{7-5}{例題8：腕が $(30°,\ 45°)$ のとき、$J^{-1}$ は1行め $(0.0647,\ 0.2415)$、2行め $(-0.3245,\ -0.3915)$ を $\det J=0.0530$ でわったもの。ほしい手の速度にかけると、関節の角速度が出る。}
\bkrecall{7-10}{同じ行列を何回もかけると、大きさのいちばん大きい固有値の向きにそろう。$(I+A\frac{t}{n})^n$ は $n$ を大きくすると $e^{At}$ に近づく。}
\bkrecall{6-11}{$\cosh x=\dfrac{e^x+e^{-x}}{2}$、$\sinh x=\dfrac{e^x-e^{-x}}{2}$。$\cosh^2x-\sinh^2x=1$。$\cosh1=1.543$、$\sinh1=1.175$。}
\end{omoidasou}

\noindent この章は\textbf{読みもの}です。例題や問題はありません。7-1〜7-10 で学んだ行列が、本物の機械やしくみの中で、どんな仕事をしているかを読んでいきます。それぞれの場面で、「どの章の、どの道具か」を（7-5）のように書いておきます。読んでいて分からなくなったら、その章にもどって確かめましょう。全部を一度に分からなくても大丈夫です。気になった場面から読んでかまいません。

2つの場面（\textbf{ロボットの逆運動学}と\textbf{長い電線}）では、少しだけ計算をします。どちらも、7-10 の「同じ計算を何回もくり返す」考えが主役です。

\begin{bkcenter}
\begin{bktable}{colspec={Q[l,m,wd=32mm] X[l,m] Q[c,m,wd=26mm]}, width=\linewidth, rowsep=1.5pt, row{1}={font=\sffamily\small}, cells={font=\small}}
場面 & 行列がしていること & 使う章 \\
ロボットの腕 & 手の位置・速さを出す。目標の位置へ近づける & 7-2〜7-7・7-10 \\
電気回路・送電線 & 電流と電圧を一度に解く。箱をつなぐ & 6-11・7-3〜7-5・7-10 \\
建物・橋・車 & 力とずれを解く。ゆれ方（モード）を見つける & 7-5・7-8・7-9 \\
画像・ゲーム & 画素の表。点を回す・動かす & 7-1・7-2・7-4・7-6・7-7 \\
検索・データ & いちばん大事な向き（固有ベクトル）をさがす & 7-8〜7-10 \\
\end{bktable}
\bktabcaption{この章で読む場面と、使う章}
\end{bkcenter}


%% ================= ロボット =================
\bksection{butsuri}{ロボットを動かす}

\bksubsection{1}{手はどこに、どんな速さで}

\label{bk:7-11:robot}4-7・7-6 の、2つの関節をもつロボットの腕を思い出しましょう。肩からひじまで $L_1=\qty{0.30}{m}$、ひじから手先まで $L_2=\qty{0.25}{m}$、肩の角 $\theta_1=30°$、ひじの角 $\theta_2=45°$ でした。

\noindent\begin{minipage}[c]{0.56\linewidth}\raggedright
\begin{itemize}[leftmargin=1.2\zw, itemsep=2pt, topsep=0pt, label={\color{\bkTT}●}]
\item \textbf{手の位置}\quad 回転行列（7-6）や同次変換行列 $A(\theta,\ L)$ の積（7-7）で、手は $(0.325,\ 0.391)$（m）。
\item \textbf{手の速さ}\quad ヤコビ行列 $J$ に関節の角速度をかける（7-2）。$J$ の1行めは $(-0.391,\ -0.241)$、2行めは $(0.325,\ 0.065)$。
\item \textbf{動けない姿勢}\quad $\det J=L_1L_2\sin\theta_2$ が0になる、腕がのびきった姿勢（特異姿勢。7-4）。
\end{itemize}
\end{minipage}\hfill
\begin{minipage}[c]{0.4\linewidth}\centering
\begin{tikzpicture}[x=1mm, y=1mm]
  \draw[bkGray, line width=0.5pt, -{Stealth[length=3pt]}] (-3,0) -- (40,0) node[anchor=west, font=\sffamily\footnotesize, text=bkInk] {$x$};
  \draw[bkGray, line width=0.5pt, -{Stealth[length=3pt]}] (0,-3) -- (0,46) node[anchor=south, font=\sffamily\footnotesize, text=bkInk] {$y$};
  \fill[bkGray!40] (-4,-2) rectangle (4,0);
  \draw[bkOrange!80!black, line width=2.6pt, line cap=round] (0,0) -- (20.78,12);
  \draw[bkBlue!80!black, line width=2.6pt, line cap=round] (20.78,12) -- (25.96,31.32);
  \filldraw[fill=white, draw=bkInk, line width=0.8pt] (0,0) circle (1.3) (20.78,12) circle (1.3);
  \fill[bkRed] (25.96,31.32) circle (1.2);
  \node[font=\sffamily\footnotesize, anchor=south, align=center] at (31,32.8) {手\\$(0.325,\ 0.391)$};
  \draw[bkInk, line width=0.5pt] (8,0) arc (0:30:8);
  \node[font=\sffamily\footnotesize, anchor=west] at (9,2.6) {$30°$};
  \node[font=\sffamily\footnotesize, anchor=north west] at (21.5,11.5) {$45°$};
  \node[font=\sffamily\footnotesize, anchor=south east] at (9.5,6.5) {$L_1$};
  \node[font=\sffamily\footnotesize, anchor=east] at (22.5,22) {$L_2$};
\end{tikzpicture}
\end{minipage}

\medskip
ここまでは「関節の角がわかっているとき、手はどこか」という計算でした。これを\textbf{順運動学}（じゅんうんどうがく）といいます。工場で本当に必要なのは、その逆です。「この部品の上に手を持っていきたい。関節を何度にすればよいか」。これを\textbf{逆運動学}（ぎゃくうんどうがく）といいます。\bkidx{ぎゃくうんどうがく}{逆運動学}

\bksubsection{2}{逆運動学：ヤコビ行列の逆行列で、少しずつ近づく}

\label{bk:7-11:gyaku}2つの関節の腕なら、4-7 で学んだように、余弦定理を使って角を一度に求められます。関節が6つある産業用ロボットの多くも、手首の3つの軸を1点で交わらせて、式で一度に解けるようにくふうしてあります。それでも、式で解けない形のロボットや、関節が7つ以上のロボット、手を道すじにそってなめらかに動かす制御では、次の「少しずつ近づける」方法が使われます。

\kviilhead{考え方：今の姿勢の $J$ で、1歩ずつ}
7-2 のヤコビ行列は、「関節の角を少し変えると、手がどれだけ動くか」を表していました。角の変化を $\Delta\vec\theta$、手の変化を $\Delta\vec p$ とすると、変化が小さいときは
\[ \Delta\vec p\approx J\,\Delta\vec\theta \]
です（$\dot{\vec p}=J\dot{\vec\theta}$ の両辺に短い時間をかけたもの）。逆に、手を $\Delta\vec p$ だけ動かしたいなら、7-5 の逆行列を使って
\kviilframe{$\Delta\vec\theta=J^{-1}\,\Delta\vec p$\\[3pt]{\small 関節の変え方 ＝ ヤコビ行列の逆行列 × 手を動かしたい量}}
とすればよいはずです。ここで、$\Delta\vec p$ は「目標の位置 − 今の手の位置」（ずれ）です。

ただし、$J$ は今の姿勢での「手の速さの向き」なので、腕が大きく回ると合わなくなります。6-2 のオイラー法が「接線の向きに1歩」だったのと同じで、1回ではぴったり目標に着きません。そこで、次の手順をくり返します。
\begin{enumerate}[label={\bkcircled[\bkTT]{\arabic*}}, leftmargin=2.4\zw, itemsep=2pt, topsep=2pt]
\item 今の角から、手の位置を計算する（順運動学。7-6・7-7）。
\item ずれ $\Delta\vec p=$ 目標 − 今の手の位置 を出す。小さければ終わり。
\item 今の角で $J$ を作り、$\Delta\vec\theta=J^{-1}\Delta\vec p$ を解く（7-5）。
\item 角を $\Delta\vec\theta$ だけ変えて、①にもどる。
\end{enumerate}
7-10 の漸化式と同じで、「今の角から次の角を作る」くり返しです。

\kviilhead{数で確かめる}
今の角 $(30°,\ 45°)$、手 $(0.325,\ 0.391)$ から、目標 $(0.20,\ 0.45)$ へ手を動かします。1回めの計算です。
\begin{bkhenkei}[0.74]
\Delta\vec p=(0.20-0.3245,\ 0.45-0.3915) & ② ずれ \\
\phantom{\Delta\vec p}=(-0.1245,\ 0.0585) & 4けたで計算 \\
0.0647\times(-0.1245)+0.2415\times0.0585=0.006073 & ③ $J^{-1}$ の1行め \\
\Delta\theta_1=0.006073\div0.0530\approx0.115 & $\det J$ でわる \\
-0.3245\times(-0.1245)-0.3915\times0.0585=0.017497 & 2行め \\
\Delta\theta_2=0.017497\div0.0530\approx0.330 & 〔rad〕 \\
\end{bkhenkei}
度に直すと、$\Delta\theta_1\approx6.6°$、$\Delta\theta_2\approx18.9°$ です（rad から度へは $\times\frac{180}{\pi}$）。7-5 例題8 では「ほしい手の速度」に $J^{-1}$ をかけましたが、ここでは「ほしい手の動き（ずれ）」にかけています。角は $(36.6°,\ 63.9°)$ になります。これを④から①へもどしてくり返すと、次の表のようになりました。

\begin{bkcenter}
\begin{bktable}{colspec={Q[c,m,wd=12mm] Q[c,m,wd=30mm] Q[c,m,wd=32mm] Q[c,m,wd=26mm]}, rowsep=1pt, row{1}={font=\sffamily\small}}
回数 & 角 $(\theta_1,\ \theta_2)$ & 手の位置（m） & ずれの長さ（m） \\
0 & $(30°,\ 45°)$ & $(0.325,\ 0.391)$ & 0.138 \\
1 & $(36.6°,\ 63.9°)$ & $(0.196,\ 0.425)$ & 0.026 \\
2 & $(41.7°,\ 54.0°)$ & $(0.199,\ 0.448)$ & 0.002 \\
3 & $(42.1°,\ 53.1°)$ & $(0.200,\ 0.450)$ & 0.00001 \\
\end{bktable}
\bktabcaption{ヤコビ行列の逆行列で、くり返し近づけた。ずれの長さは、ずれの矢印の長さ（三平方の定理）。}
\end{bkcenter}

\begin{bkfigure}
\begin{tikzpicture}[x=1.6mm, y=1.6mm]
  \draw[bkGray, line width=0.5pt, -{Stealth[length=3pt]}] (-3,0) -- (45,0) node[anchor=west, font=\sffamily\footnotesize, text=bkInk] {$x$};
  \draw[bkGray, line width=0.5pt, -{Stealth[length=3pt]}] (0,-2) -- (0,52) node[anchor=south, font=\sffamily\footnotesize, text=bkInk] {$y$};
  \fill[bkGray!40] (-3,-1.5) rectangle (3,0);
  % 0回め
  \draw[bkGray!70, line width=1.8pt, line cap=round] (0,0) -- (25.98,15) -- (32.45,39.15);
  % 1回め
  \draw[bkGreen!50, line width=1.8pt, line cap=round] (0,0) -- (24.09,17.88) -- (19.55,42.46);
  % 3回め（目標）
  \draw[bkOrange!80!black, line width=2.4pt, line cap=round] (0,0) -- (22.27,20.10);
  \draw[bkBlue!80!black, line width=2.4pt, line cap=round] (22.27,20.10) -- (20,45);
  \filldraw[fill=white, draw=bkInk, line width=0.8pt] (0,0) circle (0.8) (22.27,20.10) circle (0.8);
  \fill[bkGray!70] (32.45,39.15) circle (0.8);
  \fill[bkGreen!60!black] (19.55,42.46) circle (0.8);
  \draw[bkRed, line width=1pt] (20,45) circle (1.4);
  \draw[bkRed, line width=0.8pt] (18.2,45) -- (21.8,45) (20,43.2) -- (20,46.8);
  \node[font=\sffamily\footnotesize, anchor=south, text=bkRed!80!black] at (20,47) {目標 $(0.20,\ 0.45)$};
  \node[font=\sffamily\footnotesize, anchor=west, text=bkGray] at (33.6,39.2) {0回め};
  \node[font=\sffamily\footnotesize, anchor=east, text=bkGreen!50!black] at (18.2,41.3) {1回め};
  \node[font=\sffamily\footnotesize, anchor=north west, text=bkOrange!80!black] at (23,19.6) {3回め};
  \draw[bkInk, line width=0.5pt, dash pattern=on 1.5pt off 1pt, -{Stealth[length=2.5pt]}] (31.6,39.6) -- (20.6,42.3);
\end{tikzpicture}
\bkfigcaption{くり返すたびに、手（先の点）が目標（赤い十字）に近づく。灰色が0回め、緑が1回め、オレンジと青が3回め（ほぼ目標）。}
\label{bk:7-11:fig-ik}
\end{bkfigure}

1回めで、ずれは0.138\,m から0.026\,m に、2回めで0.002\,m に、3回めで0.00001\,m（0.01\,mm）になりました。目標に近づくほど、$J$ の「接線」がよく合うので、近づき方がどんどん速くなります。4-7 の余弦定理で求めた答え（$\theta_1=42.1°$、$\theta_2=53.1°$）とも一致します。

関節が6つでも、手順はまったく同じです。$J$ が $6\times6$ の行列になり、③で6つの文字の連立方程式を解くだけです（7-5 の消去法）。ロボットの制御の計算機は、この計算を1秒に何百〜何千回もくり返しながら、手を目標の道すじにそって動かしています。

\Needspace{10\baselineskip}
\begin{seikaku}
このくり返しは、大学の本では\textbf{ニュートン・ラフソン法}（多変数のニュートン法）の仲間として出てきます。うまくいくのは、目標がとどく所にあり、とちゅうで $\det J=0$ の特異姿勢（7-4）を通らないときです。特異姿勢の近くでは $J^{-1}$ の数がとても大きくなり、関節を急に大きく回す答えが出てしまいます。実際のロボットでは、1回に変える角に上限をつけたり、特異姿勢をさける道すじを選んだりします。$J$ の作り方（偏微分）は、第8巻の 8-3 で学びます。
\end{seikaku}

\bksubsection{3}{ドローン・人工衛星・VR ゴーグルの向き}

\noindent\begin{minipage}[c]{0.55\linewidth}\raggedright
\label{bk:7-11:muki}ドローンや人工衛星は、ジャイロセンサーで「機体がどの向きに、どれだけの速さで回っているか」を1秒に何百〜何千回もはかり、3次元の回転行列（7-7）で機体の向きを更新しつづけます。ロール・ピッチ・ヨーの3つの角だけで向きを表すと、ピッチが $\pm90°$ になるとジンバルロック（7-7）が起きるので、宇宙船やドローンでは、4つの数で向きを表すクォータニオン（6-9）もよく使われます。
\end{minipage}\hfill
\begin{minipage}[c]{0.42\linewidth}\centering
\begin{tikzpicture}[x=0.75mm, y=0.75mm]
  % 機体（ななめから見たX形）
  \draw[bkInk, line width=1.4pt] (-14,-5) -- (14,5) (-14,5) -- (14,-5);
  \foreach \p in {(-14,-5),(14,5),(-14,5),(14,-5)} {\draw[bkBlue!70!black, line width=0.9pt] \p ellipse (6 and 1.6);}
  \filldraw[fill=bkGray!60, draw=bkInk, line width=0.6pt] (-3,-2) rectangle (3,2);
  % 3本の軸
  \draw[bk arrow, bkRed, line width=1pt] (0,0) -- (26,-9);
  \node[font=\sffamily\footnotesize, text=bkRed!80!black, anchor=north, align=center] at (24,-10) {ロール\\（前後の軸）};
  \draw[bk arrow, bkGreen!60!black, line width=1pt] (0,0) -- (-22,-10);
  \node[font=\sffamily\footnotesize, text=bkGreen!40!black, anchor=north, align=center] at (-20,-11.5) {ピッチ\\（左右の軸）};
  \draw[bk arrow, bkOrange!80!black, line width=1pt] (0,0) -- (0,20);
  \node[font=\sffamily\footnotesize, text=bkOrange!70!black, anchor=south] at (0,20.5) {ヨー（上下の軸）};
\end{tikzpicture}
\end{minipage}

\medskip
VR ゴーグルや3次元のゲームでは、物の点を「世界の座標」に置く $4\times4$ の行列と、カメラ（頭）から見た座標に直す行列をかけて、画面を作ります（7-7）。何万もの点に、1秒に何十回も同じ行列をかけるので、そのための専用の計算回路（GPU）が使われています。電卓が $\sin$・$\cos$ を計算するときにも、回転をくり返す CORDIC という方法が使われています（7-6 のコラム）。

%% ================= 電気回路 =================
\bksection{butsuri}{電気回路と送電線}

\bksubsection{4}{ばねの網と、抵抗の網}

\label{bk:7-11:ami}7-2 の物理・工学ではこう使うで、かべにばね2本（ばね定数\qty{100}{N/m}）を直列につないだ「ばねの網」を見ました。点1・点2のずれ $u_1$、$u_2$〔m〕と、そこにかける力 $F_1$、$F_2$〔N〕の関係は、
\[ \begin{pmatrix}F_1\\F_2\end{pmatrix}=\begin{pmatrix}200&-100\\-100&100\end{pmatrix}\begin{pmatrix}u_1\\u_2\end{pmatrix} \]
でした。こんどは、電気で同じ形を作ります。

\begin{bkfigure}
\begin{tikzpicture}[x=1mm, y=1mm]
  % ばねの網
  \begin{scope}
    \fill[bkGray!35] (-4,-6) rectangle (0,8);
    \draw[bkInk, line width=0.9pt] (0,-6) -- (0,8);
    \kviilspring{0}{20}{1}
    \filldraw[fill=bkOrange!25, draw=bkOrange!80!black, line width=0.8pt] (20,-3) rectangle (26,5);
    \kviilspring{26}{46}{1}
    \filldraw[fill=bkBlue!20, draw=bkBlue!80!black, line width=0.8pt] (46,-3) rectangle (52,5);
    \node[font=\sffamily\footnotesize] at (23,1) {1};
    \node[font=\sffamily\footnotesize] at (49,1) {2};
    \node[font=\sffamily\footnotesize, anchor=south] at (10,4.5) {$k$};
    \node[font=\sffamily\footnotesize, anchor=south] at (36,4.5) {$k$};
    \draw[bk arrow, bkRed, line width=0.9pt] (20,8.5) -- (27,8.5);
    \node[font=\sffamily\footnotesize, anchor=south, text=bkRed!80!black] at (23.5,9.5) {$F_1$};
    \draw[bk arrow, bkRed, line width=0.9pt] (52.5,1) -- (59,1);
    \node[font=\sffamily\footnotesize, anchor=west, text=bkRed!80!black] at (59.5,1) {$F_2$};
    \node[font=\sffamily\small, anchor=north] at (28,-17) {ばねの網：ずれ $u$ と力 $F$};
  \end{scope}
  % 抵抗の網
  \begin{scope}[xshift=72mm]
    \draw[bkInk, line width=0.8pt] (0,8) -- (6,8) (14,8) -- (26,8) (34,8) -- (44,8);
    \kviilres{10}{8}
    \kviilres{30}{8}
    \draw[bkInk, line width=0.8pt] (0,-6) -- (44,-6);
    \draw[bkInk, line width=0.8pt] (0,8) -- (0,-6);
    \fill[bkInk] (20,8) circle (1) (44,8) circle (1);
    \node[font=\sffamily\footnotesize, anchor=south] at (20,9.5) {1};
    \node[font=\sffamily\footnotesize, anchor=south] at (44,9.5) {2};
    \node[font=\sffamily\footnotesize, anchor=north] at (10,5.5) {$R$};
    \node[font=\sffamily\footnotesize, anchor=north] at (30,5.5) {$R$};
    \node[font=\sffamily\footnotesize, anchor=north] at (22,-7) {電位0（地面）};
    \draw[bk arrow, bkRed, line width=0.9pt] (20,0) -- (20,6.5);
    \node[font=\sffamily\footnotesize, anchor=north, text=bkRed!80!black] at (20,-0.5) {$I_1$};
    \draw[bk arrow, bkRed, line width=0.9pt] (44,0) -- (44,6.5);
    \node[font=\sffamily\footnotesize, anchor=north, text=bkRed!80!black] at (44,-0.5) {$I_2$};
    \node[font=\sffamily\small, anchor=north] at (22,-17) {抵抗の網：電位 $V$ と電流 $I$};
  \end{scope}
\end{tikzpicture}
\bkfigcaption{左：かべとばね2本。右：電位0の線（地面）と抵抗2本。点1・点2に、外から電流 $I_1$、$I_2$ を流しこむ。}
\label{bk:7-11:fig-ami}
\end{bkfigure}

抵抗 $R$〔Ω〕の電気の通りやすさを $G=\dfrac1R$（単位 S、ジーメンス）と書きます。オームの法則は「電流＝$G$×電圧」です。点1・点2の電位（地面からの電圧）を $V_1$、$V_2$〔V〕とし、外から流しこむ電流を $I_1$、$I_2$〔A〕とします。流しこんだ電流は、つながった抵抗を通って出ていきます（3-6 のきまり1：合わさった電流はたし算）。
\begin{bkhenkei}[0.48]
I_1=GV_1+G(V_1-V_2) & 点1から地面へと、点1から点2へ \\
I_2=G(V_2-V_1) & 点2から点1へ \\
\end{bkhenkei}
同じ文字でまとめて、係数を行列にならべると（7-1）、
\[ \begin{pmatrix}I_1\\I_2\end{pmatrix}=\begin{pmatrix}2G&-G\\-G&G\end{pmatrix}\begin{pmatrix}V_1\\V_2\end{pmatrix} \]
ばねの網と\textbf{まったく同じ形の行列}です。$G=\qty{0.1}{S}$（$R=\qty{10}{\ohm}$）、$V_1=\qty{1}{V}$、$V_2=\qty{3}{V}$ なら、
\[ I_1=0.2\times1-0.1\times3=\qty{-0.1}{A},\qquad I_2=-0.1\times1+0.1\times3=\qty{0.2}{A} \]7-2 のばねの $F_1=-1$、$F_2=2$ と、数の並びまで同じです。

\begin{bkcenter}
\begin{bktable}{colspec={Q[c,m,wd=40mm] Q[c,m,wd=40mm]}, rowsep=1pt, row{1}={font=\sffamily\small}}
ばねの網 & 抵抗の網 \\
ずれ $u$〔m〕 & 電位 $V$〔V〕 \\
力 $F$〔N〕 & 流しこむ電流 $I$〔A〕 \\
ばね定数 $k$〔N/m〕 & 通りやすさ $G=\frac1R$〔S〕 \\
\end{bktable}
\bktabcaption{ばねの網と抵抗の網の対応。6-15 の相似の表（力↔電圧）とは組み方がちがう、もう1つの相似（力↔電流）。}
\end{bkcenter}
6-15 の相似の表では、力を電圧に対応させました。ここでは力を電流に対応させています。どちらも正しい相似で、問題によって使いやすいほうを選びます。網の形（どの点とどの点がつながっているか）が同じなら、行列も同じ形になるのです。電子回路のシミュレーター（7-5）は、何万もの点の電位を、この形の連立方程式 $K\vec V=\vec I$ で一度に求めています（物理 電-5）。トランジスタのように電流と電圧が比例しない部品があるときは、②と同じく、少しずつ直す連立方程式をくり返し解きます（物理 原-10）。

\bksubsection{5}{箱をつなぐ：四端子定数}

\label{bk:7-11:hako}7-3 では、入口2本・出口2本の線の「箱」を、2×2 の行列で表しました。出口の電圧・電流の組を入れると、入口の組が出てくる行列です。
\[ \text{直列の抵抗の箱}\ \begin{pmatrix}1&R\\0&1\end{pmatrix},\qquad \text{並列の抵抗の箱}\ \begin{pmatrix}1&0\\\dfrac1R&1\end{pmatrix} \]
箱をつなぐと、行列のかけ算になりました（入口側の箱を左に書く）。電気の本では、この4つの数を $A$、$B$、$C$、$D$ と書いて\textbf{四端子定数}とよび、抵抗・コイル・コンデンサーだけの箱では、いつも $AD-BC=1$（7-4 の行列式が1）になります。音を整える回路、雑音を取りのぞくフィルター、アンテナと機械をつなぐ回路などは、箱をいくつもつないで設計します。カメラのレンズの光線の行列（7-3・物理 波-12）も、同じ形の行列のかけ算でした。

\bksubsection{6}{長い電線は、小さな箱のくさり：分布定数回路}

\label{bk:7-11:bunpu}数百\,km の送電線や、テレビのアンテナ線のように長い電線では、線そのものが抵抗やコイル・コンデンサーのはたらきを、全体にうすく広がった形でもっています。このような回路を\textbf{分布定数回路}（ぶんぷていすうかいろ）といいます。\bkidx{ぶんぷていすうかいろ}{分布定数回路}部品が1か所にまとまった、これまでの回路とのちがいです。

\begin{bkfigure}
\begin{tikzpicture}[x=1mm, y=1mm]
  \begin{scope}[xscale=0.88]
  \draw[bkInk, line width=0.8pt] (0,10) -- (2,10) (0,-2) -- (122,-2);
  \foreach \k in {0,...,5} {
    \begin{scope}[xshift={20*\k mm}]
      \draw[bkInk, line width=0.8pt] (2,10) -- (4,10) (12,10) -- (22,10);
      \kviilres{8}{10}
      \kviilresv{17}{4}
      \draw[bkInk, line width=0.8pt] (17,10) -- (17,8) (17,0) -- (17,-2);
      \draw[bkTheme, line width=0.6pt, dash pattern=on 1.2pt off 1pt, rounded corners=1mm] (2.5,-4) rectangle (20.5,14);
    \end{scope}}
  \fill[bkInk] (0,10) circle (0.9) (0,-2) circle (0.9) (122,10) circle (0.9) (122,-2) circle (0.9);
  \node[font=\sffamily\footnotesize, anchor=east] at (-1.5,4) {入口};
  \node[font=\sffamily\footnotesize, anchor=west] at (123.5,4) {出口};
  \node[font=\sffamily\footnotesize, anchor=south] at (11.5,14.5) {長さ $\Delta x$ の箱};
  \end{scope}
\end{tikzpicture}
\bkfigcaption{長い電線を、長さ $\Delta x$ の小さな箱のくさりと見る。1つの箱は、上の線の抵抗（直列）と、2本の線のあいだのもれ（並列）でできている。}
\label{bk:7-11:fig-kusari}
\end{bkfigure}

\kviilhead{数を簡単にした長い線で計算する}
しくみを見るために、数を簡単にした線を考えます。上の線の抵抗が1\,km あたり $r=\qty{1}{\ohm}$、2本の線のあいだのもれ（電気の通りやすさ）が1\,km あたり $g=\qty{1}{S}$、線の長さが1\,km とします。これを $n$ 個の小さな箱に分けると、1つの箱の長さは $\Delta x=\frac1n$\,km で、その箱は「直列の抵抗 $r\Delta x$」と「並列の通りやすさ $g\Delta x$」をつないだものです（7-3）。
\[ \begin{pmatrix}1&r\Delta x\\0&1\end{pmatrix}\begin{pmatrix}1&0\\g\Delta x&1\end{pmatrix}=\begin{pmatrix}1+rg\Delta x^2&r\Delta x\\g\Delta x&1\end{pmatrix} \]
線全体は、この小さな箱を $n$ 個つないだものなので、この行列の $n$ 乗です（7-3・7-10）。Python で計算すると、次のようになりました（小数第3位まで）。

\begin{bkcenter}
\begin{bktable}{colspec={Q[c,m,wd=20mm] Q[c,m,wd=20mm] Q[c,m,wd=20mm] Q[c,m,wd=20mm] Q[c,m,wd=18mm]}, rowsep=1pt, row{1}={font=\sffamily\small}}
箱の数 $n$ & 左上 & 右上 & 左下 & 行列式 \\
10 & 1.601 & 1.173 & 1.173 & 1 \\
100 & 1.549 & 1.175 & 1.175 & 1 \\
1000 & 1.544 & 1.175 & 1.175 & 1 \\
限りなく多く & 1.543 & 1.175 & 1.175 & 1 \\
\end{bktable}
\bktabcaption{小さな箱を $n$ 個つないだ行列の成分（右下は左上とほぼ同じ数に近づく）。}
\end{bkcenter}
箱をこまかくしていくと、成分は1.543と1.175に近づきます。この2つの数に見覚えはありませんか。6-11 の $\cosh1=1.543$ と $\sinh1=1.175$ です。行列式は、どの $n$ でもぴったり1（7-4：箱の行列の行列式は1で、積の行列式は行列式の積）。限りなく細かくした先の行列でも、$\cosh^21-\sinh^21=1$（6-11）です。

\kviilhead{専門書の式}
一般の長い線では、1\,m あたりの直列の部分を $z$、並列の部分を $y$ として、長さ $\ell$ の線の行列は
\kviilframe{$\begin{pmatrix}\cosh\gamma\ell&Z_0\sinh\gamma\ell\\[3pt]\dfrac{1}{Z_0}\sinh\gamma\ell&\cosh\gamma\ell\end{pmatrix}$\qquad $\gamma=\sqrt{zy}$、\ $Z_0=\sqrt{\dfrac zy}$}
になります。上の計算は $z=r=1$、$y=g=1$、$\ell=1$ で、$\gamma=1$、$Z_0=1$ の場合でした。$\gamma$ を\textbf{伝搬定数}（でんぱんていすう）、$Z_0$ を\textbf{特性インピーダンス}（とくせいインピーダンス）といいます。\bkidx{とくせいいんぴーだんす}{特性インピーダンス}テレビのアンテナ線（同軸ケーブル）に「75\,Ω」と書いてあるのは、この $Z_0$ のことです。

本物の送電線では、$z$ に 6-10 のコイルの $j\omega L$、$y$ にコンデンサーの $j\omega C$ が入るので、$\gamma$ は複素数になります。すると $\cosh$・$\sinh$ が、6-10 の $e^{j\theta}$ のように「回る」形になり、電線の上を電圧の波が伝わっていくようすが表れます。くわしくは、物理シリーズの電磁気学（電-13）で学びます。

\begin{seikaku}
「小さな箱を $n$ 個つないで、$n$ を大きくする」は、7-10 のもっと知りたいの「$\left(I+A\frac{t}{n}\right)^n$ が $e^{At}$ に近づく」と同じ形です。小さな箱の行列は、$\Delta x^2$ の項を小さいとして省くと $I+M\Delta x$（$M$ は1行め $(0,\ z)$、2行め $(y,\ 0)$ の行列）と書けるので、長さ $\ell$ の線の行列は $e^{M\ell}$ です。$M$ の固有値は $\pm\sqrt{zy}=\pm\gamma$ で（7-8）、$e^{\gamma\ell}$ と $e^{-\gamma\ell}$ の平均と差の半分が $\cosh\gamma\ell$ と $\sinh\gamma\ell$ です（6-11）。行列の指数関数 $e^{M\ell}$ は、第9巻（9-11）でくわしく学びます。
\end{seikaku}

%% ================= 建物・画像・検索 =================
\bksection{butsuri}{建物・画像・データ}

\bksubsection{7}{建物と橋：力とずれ、ゆれ方}

\label{bk:7-11:tatemono}橋やビルを設計するときは、建物を小さな部品に分けて、つながる点ごとに力のつり合いの式を立てます（有限要素法。7-5）。④のばねの網を何百万点にも広げたもので、$K\vec u=\vec F$ の形の連立方程式を、7-5 の消去法の考えを使ったプログラムで解きます。$K$ は④と同じく、ほとんどの成分が0の大きな行列です（となりの点としかつながっていないから）。0の所を計算しない工夫で、何百万元の式も解けるのです。

\begin{bkfigure}
\begin{tikzpicture}[x=1mm, y=1mm]
  \foreach \xs/\da/\db/\nm in {0/0/0/{ゆれていない}, 44/3/5.5/{1次モード：同じ向き}, 88/2.5/-3/{2次モード：反対向き}} {
    \begin{scope}[xshift=\xs mm]
      \fill[bkGray!40] (-4,-2) rectangle (28,0);
      \draw[bkInk, line width=0.6pt] (-4,0) -- (28,0);
      % 柱（1階）
      \draw[bkBlue!70!black, line width=1.2pt] (2,0) .. controls (2,7) and ({2+\da},7) .. ({2+\da},14);
      \draw[bkBlue!70!black, line width=1.2pt] (22,0) .. controls (22,7) and ({22+\da},7) .. ({22+\da},14);
      % 2階の柱
      \draw[bkBlue!70!black, line width=1.2pt] ({2+\da},14) .. controls ({2+\da},21) and ({2+\db},21) .. ({2+\db},28);
      \draw[bkBlue!70!black, line width=1.2pt] ({22+\da},14) .. controls ({22+\da},21) and ({22+\db},21) .. ({22+\db},28);
      % 床
      \filldraw[fill=bkOrange!35, draw=bkOrange!80!black, line width=0.8pt] ({0+\da},13) rectangle ({24+\da},15.5);
      \filldraw[fill=bkOrange!35, draw=bkOrange!80!black, line width=0.8pt] ({0+\db},27) rectangle ({24+\db},29.5);
      \node[font=\sffamily\footnotesize, anchor=north, align=center] at (12,-3) {\nm};
    \end{scope}}
  \draw[bk arrow, bkRed, line width=0.9pt] (70,32) -- (75,32);
  \draw[bk arrow, bkRed, line width=0.9pt] (69,18.5) -- (72,18.5);
  \draw[bk arrow, bkRed, line width=0.9pt] (114,18.5) -- (117,18.5);
  \draw[bk arrow, bkRed, line width=0.9pt] (110,32) -- (106,32);
\end{tikzpicture}
\bkfigcaption{2階建ての建物の2つの振動モード（固有ベクトル）。2つの床が同じ向きにゆれる形と、反対向きにゆれる形。7-8 の台車2つの「いっしょ」「反対」と同じ。}
\end{bkfigure}


地震や風でのゆれ方は、7-8・7-9 の固有値と固有ベクトルで調べます。各階の重さと、階と階をつなぐ柱の強さから行列を作ると、固有ベクトルが「ゆれ方の形（モード）」、固有値が「ゆれやすい振動数」を教えてくれます。ロンドンのミレニアム・ブリッジが人の歩くリズムでゆれた話（7-8）、車のバウンスとピッチ（7-9）、タイヤのバランス（7-8）も、同じ道具です。機械の部品の主応力（7-8）は、ひずみゲージではかったのびから応力の行列を作り、その固有値として求めます。

\bksubsection{8}{画像とコンピュータグラフィックス}

\label{bk:7-11:gazou}写真は、画素の明るさの数がたてと横に並んだ行列です（7-1）。全体を明るくするのは実数倍、2まいを重ねるのは和でした。写真を回すアプリは、画素の1つ1つに回転行列をかけ（7-6）、ゲームは何万もの点に $4\times4$ の行列をかけます（7-7）。見えない裏側の三角形は、2×2 の行列式の符号で見分けてかかないようにしています（7-4）。


\noindent\begin{minipage}[c]{0.62\linewidth}\raggedright
写真を小さなファイルにするときにも、行列が働いています。ふつうの写真（JPEG という形式）は、画像を $8\times8$ 画素の小さなブロックに分け、それぞれのブロックを「なめらかな部分」と「細かいもようの部分」に分けて表し直し、目で気づきにくい細かい部分の数を少なくして記録します。この「表し直し」は、$8\times8$ の行列のかけ算です（9-12 で読むフーリエの考えの親せき）。また、大きな行列を、いちばん大事な向きの組み合わせだけで近似する\textbf{特異値分解}という方法もあり、画像の圧縮やデータの整理に使われます（7-9 の対角化の親せき）。
\end{minipage}\hfill
\begin{minipage}[c]{0.34\linewidth}\centering
\begin{tikzpicture}[x=1mm, y=1mm]
  \foreach \i/\j/\v in {0/0/20,1/0/60,2/0/140,3/0/200, 0/1/40,1/1/120,2/1/200,3/1/240, 0/2/90,1/2/180,2/2/230,3/2/250, 0/3/150,1/3/220,2/3/250,3/3/255} {
    \pgfmathsetmacro\g{100-\v/255*100}
    \fill[black!\g] ({8*\i},{-8*\j}) rectangle ++(8,-8);
    \pgfmathsetmacro\tc{\v>128?0:1}
    \ifdim\tc pt>0.5pt \node[font=\sffamily\footnotesize, text=white] at ({8*\i+4},{-8*\j-4}) {\v};\else\node[font=\sffamily\footnotesize, text=bkInk] at ({8*\i+4},{-8*\j-4}) {\v};\fi}
  \draw[bkInk, line width=0.6pt] (0,0) rectangle (32,-32);
\end{tikzpicture}\par\vspace{1pt}
{\footnotesize 4×4 画素の明るさ（0〜255）}
\end{minipage}


\bksubsection{9}{検索エンジンの順位：固有ベクトル}

\label{bk:7-11:kensaku}インターネットで検索すると、たくさんのページが順位をつけて表示されます。1998年ころ、アメリカの大学院生だったブリンとペイジは、「大事なページからたくさんリンクされているページは大事」という考えで順位を決める方法を発表しました。\textbf{ページランク}とよばれる方法です。\bkidx{ぺーじらんく}{ページランク}その中心は、7-10 の天気の行列と同じ「移り方の表」と、その固有ベクトルです。

\noindent\begin{minipage}[c]{0.4\linewidth}\centering
\begin{tikzpicture}[x=1mm, y=1mm]
  \kviilpage{(0,0)}{A}{bkOrange}
  \kviilpage{(34,0)}{B}{bkGreen}
  \kviilpage{(17,-30)}{C}{bkBlue}
  \draw[bk link, bkInk, line width=0.9pt] (8.5,2) -- (25.5,2);
  \draw[bk link, bkInk, line width=0.9pt] (3,-9.5) -- (10.5,-21);
  \draw[bk link, bkInk, line width=0.9pt] (31,-9.5) -- (23.5,-21);
  \draw[bk link, bkInk, line width=0.9pt] (9.5,-25) .. controls (-6,-22) and (-8,-14) .. (-4,-9.5);
\end{tikzpicture}
\end{minipage}\hfill
\begin{minipage}[c]{0.56\linewidth}\raggedright
3つのページ A・B・C のリンク（矢印）。A は B と C へ、B は C へ、C は A へリンクしている。
\end{minipage}

\medskip
ページを読む人が、リンクを1つ選んでとなりのページへ移ることをくり返すとします。A にいる人は、B か C へ半分ずつ移ります。B にいる人は C へ、C にいる人は A へ移ります。今いる確率の組 $(a,\ b,\ c)$ は、7-10② の天気と同じように、移り方の表（各列の和が1）をかけて次の組になります。
\[ \begin{pmatrix}a_{n+1}\\b_{n+1}\\c_{n+1}\end{pmatrix}=\begin{pmatrix}0&0&1\\0.5&0&0\\0.5&1&0\end{pmatrix}\begin{pmatrix}a_n\\b_n\\c_n\end{pmatrix} \]
はじめ3つに $\frac13$ ずついたとして、くり返しかけると（7-8④ のべき乗法）、$(0.333,\ 0.167,\ 0.5)$、$(0.5,\ 0.167,\ 0.333)$、$(0.333,\ 0.25,\ 0.417)$、…と、ゆれながら $(0.4,\ 0.2,\ 0.4)$ に近づきます。この組を行列にかけると、$(0.4,\ 0.5\times0.4,\ 0.5\times0.4+0.2)=(0.4,\ 0.2,\ 0.4)$ と変わりません。\textbf{固有値1の固有ベクトル}です。A と C が同じくらい大事で、B はその半分、という順位がつきました。


本物の検索エンジンでは、ページは何十億もあります。それでも、するのは「大きな行列を何回もかけて、固有値1の固有ベクトルを見つける」ことです（実際には、リンクと関係なくほかのページへ移る分もまぜて、くり返しが必ず落ち着くようにします）。7-10 のくり返しの考えが、そのまま使われています。
\begin{bkfigure}
\begin{tikzpicture}
\begin{axis}[bkplot, axis x line=bottom, axis y line=left, width=112mm, height=44mm, xmin=0, xmax=12.8, ymin=0, ymax=0.6, xtick={0,2,4,6,8,10,12}, ytick={0.2,0.4}, yticklabel style={/pgf/number format/fixed}, xlabel={かけた回数 $n$}, ylabel={確率}]
  \addplot[bkthin, dashed, domain=0:12.5] {0.4};
  \addplot[bkthin, dashed, domain=0:12.5] {0.2};
  \addplot[bkcurve2, line width=1.1pt, mark=*, mark size=1.5pt] coordinates {(0,0.333) (1,0.333) (2,0.5) (3,0.333) (4,0.417) (5,0.417) (6,0.375) (7,0.417) (8,0.396) (9,0.396) (10,0.406) (11,0.396) (12,0.401)};
  \addplot[bkGreen!60!black, line width=1.1pt, mark=square*, mark size=1.4pt] coordinates {(0,0.333) (1,0.167) (2,0.167) (3,0.25) (4,0.167) (5,0.208) (6,0.208) (7,0.188) (8,0.208) (9,0.198) (10,0.198) (11,0.203) (12,0.198)};
  \addplot[bkcurve, line width=1.1pt, dashed, mark=triangle*, mark size=2pt] coordinates {(0,0.333) (1,0.5) (2,0.333) (3,0.417) (4,0.417) (5,0.375) (6,0.417) (7,0.396) (8,0.396) (9,0.406) (10,0.396) (11,0.401) (12,0.401)};
  \node[font=\sffamily\footnotesize, text=bkOrange!80!black, anchor=south west, fill=white, inner sep=1pt] at (axis cs:2.2,0.51) {A};
  \node[font=\sffamily\footnotesize, text=bkBlue, anchor=south east, fill=white, inner sep=1pt] at (axis cs:0.9,0.51) {C};
  \node[font=\sffamily\footnotesize, text=bkGreen!50!black, anchor=north, fill=white, inner sep=1pt] at (axis cs:1.5,0.15) {B};
\end{axis}
\end{tikzpicture}
\bkfigcaption{A（オレンジ）・B（緑）・C（青の点線）にいる確率。ゆれながら、0.4・0.2・0.4（灰色の点線）に落ち着く。固有値1以外の固有値は複素数（大きさ約0.71）なので、回りながら（ゆれながら）消えていく（7-10⑦）。}
\end{bkfigure}

\bksubsection{10}{位置を決める：GPS とセンサー}

\label{bk:7-11:gps}スマートフォンの GPS は、高さ約\qty{20000}{km}の空を回る人工衛星から届く電波の時刻から、衛星までの距離を求めます。わからない数は、自分の位置 $x$、$y$、$z$ と、スマートフォンの時計のずれの4つです。時計が100万分の1秒ずれるだけで、光の速さ（1秒に約30万\,km）では\qty{300}{m}のずれになるので、時計のずれも答えの1つとして求めます。だから、4つ以上の衛星が必要です。
\begin{bkfigure}
\begin{tikzpicture}[x=1mm, y=1mm]
  \fill[bkGreen!15] (-60,-6) .. controls (-20,4) and (20,4) .. (60,-6) -- (60,-12) -- (-60,-12) -- cycle;
  \draw[bkGreen!50!black, line width=0.8pt] (-60,-6) .. controls (-20,4) and (20,4) .. (60,-6);
  \node[font=\sffamily\footnotesize, anchor=north] at (40,-6) {地面};
  \coordinate (P) at (0,1.5);
  \foreach \x/\y/\n in {-45/38/1, -15/48/2, 18/46/3, 46/34/4} {
    \draw[bkGray, line width=0.6pt, dash pattern=on 2pt off 1.5pt] (\x,\y) -- (P);
    \begin{scope}[shift={(\x,\y)}]
      \fill[bkInk!70] (-2.2,-1.4) rectangle (2.2,1.4);
      \fill[bkBlue!50] (-7,-0.9) rectangle (-2.6,0.9) (2.6,-0.9) rectangle (7,0.9);
    \end{scope}
    \node[font=\sffamily\footnotesize, anchor=south] at (\x,\y+2) {衛星\n};}
  \draw[bkRed!70, line width=0.7pt] (P) ++(110:6) arc (110:70:6);
  \fill[bkRed] (P) circle (1.3);
  \node[font=\sffamily\footnotesize, anchor=north, text=bkRed!80!black, fill=bkGreen!15, inner sep=1pt] at (0,-1) {スマートフォン $(x,\ y,\ z)$};
  \node[font=\sffamily\footnotesize, anchor=west, fill=white, inner sep=1pt] at (-34,22) {$r_1$};
  \node[font=\sffamily\footnotesize, anchor=west, fill=white, inner sep=1pt] at (-6.5,26) {$r_2$};
  \node[font=\sffamily\footnotesize, anchor=west, fill=white, inner sep=1pt] at (10,26) {$r_3$};
  \node[font=\sffamily\footnotesize, anchor=west, fill=white, inner sep=1pt] at (27,19) {$r_4$};
\end{tikzpicture}
\bkfigcaption{4つの衛星までの距離 $r_1$〜$r_4$（点線）。わからない数は、位置 $x$、$y$、$z$ と時計のずれの4つなので、4つの距離の式を連立させて求める。}
\end{bkfigure}

距離の式には2乗が入るので、そのままでは連立1次方程式になりません。そこで、②の逆運動学と同じように、「近い位置を仮に決める → そこからの少しのずれについての連立1次方程式を解く（7-5）→ 位置を直す」をくり返します。5つ以上の衛星が見えるときは、式の数がわからない数より多いので、7-5 の最小二乗法で、いちばんよく合う位置を求めます。

工場の温度センサーや圧力センサーも、はかった電圧から温度などを計算するための「目もりの直線」を、最小二乗法で決めておきます（センサーの校正。7-5）。ロボットのセンサーも、ビルのひずみゲージも、こうして目もりを合わせてから使われています。

\bksubsection{11}{データと人工知能}

\label{bk:7-11:data}クラス全員の身長と体重のように、たくさんの人について、いくつもの量をはかったデータは、人を行・量を列とする大きな行列です。7-9 で見たように、ばらつきを表す対称な行列を対角化すると、データの「いちばん大事な向き」が固有ベクトルとして見つかります（主成分分析）。顔の写真の見分けや、工場の機械の調子の見守りに使われています。
\begin{bkcenter}
\begin{bktable}{colspec={Q[c,m,wd=18mm] Q[c,m,wd=22mm] Q[c,m,wd=26mm] Q[c,m,wd=20mm]}, rowsep=1pt, row{1}={font=\sffamily\small}}
 & 身長（cm） & 腕の長さ（cm） & 体重（kg） \\
1人め & 152 & 66 & 44 \\
2人め & 160 & 70 & 50 \\
3人め & 147 & 63 & 41 \\
$\vdots$ & $\vdots$ & $\vdots$ & $\vdots$ \\
\end{bktable}
\bktabcaption{データの行列の例（数は説明のためのもの）。人が行、はかった量が列。何百人・何十の量でも、1つの行列にまとめて計算する。}
\end{bkcenter}

写真に写っている物を見分けたり、文章を作ったりする人工知能も、その計算の大部分は、たくさんの数の組に行列をかけて、たし合わせることのくり返しです。何億もの数が入った行列をすばやくかけるために、ゲームの画面を作るのと同じ種類の計算回路（GPU）が使われています。

\bksubsection{12}{制御：第9巻へ}

\label{bk:7-11:seigyo}7-8 で、ばねやダンパーの動きを $\dot{\vec x}=A\vec x$ と書き、固有値が 6-12 の $s$ と同じものだと学びました。機械を思いどおりに動かすときは、これにモーターなどで加える力 $\vec u$ を足した
\[ \dot{\vec x}=A\vec x+B\vec u \]
の形を使います。これを\textbf{状態方程式}といいます。
\begin{bkfigure}
\begin{tikzpicture}[x=1mm, y=1mm]
  \node[bk node, minimum width=24mm, align=center, font=\sffamily\footnotesize] (c) at (0,0) {コンピュータ\\（行列の計算）};
  \node[bk node, minimum width=20mm, align=center, font=\sffamily\footnotesize] (m) at (38,0) {モーター};
  \node[bk node, minimum width=24mm, align=center, font=\sffamily\footnotesize] (k) at (76,0) {機械\\（状態 $\vec x$）};
  \node[bk node, minimum width=20mm, align=center, font=\sffamily\footnotesize] (s) at (38,-16) {センサー};
  \node[font=\sffamily\footnotesize, anchor=east] (g) at (-20,0) {目標};
  \draw[bk link] (g) -- (c);
  \draw[bk link] (c) -- node[bk link label, above] {$\vec u$} (m);
  \draw[bk link] (m) -- (k);
  \draw[bk link] (k.south) |- (s);
  \draw[bk link] (s) -| (c.south);
\end{tikzpicture}
\bkfigcaption{機械を思いどおりに動かすしくみ（制御）。センサーではかった状態から、コンピュータが行列の計算で力 $\vec u$ を決め、モーターで機械を動かす。これを短い時間ごとにくり返す。}
\end{bkfigure}ドローンの姿勢、エレベーターの位置、エアコンの温度などを、この式で表し、固有値が「ゆれずに・こわれずに落ち着く」場所にくるように $\vec u$ の決め方を選びます。7-10 で見たように、コンピュータの中のくり返しでは、固有値の大きさが1より小さいことが条件でした。これらは、第9巻（9-10・9-11）の主役です。

%% ================= 専門書 =================
\begin{senmonsho}[線形代数の本の書き方]
\begin{bkhonbun}
$m\times n$ 行列 $\symbf{A}$ と $\symbf{x}\in\mathbb{R}^n$ に対し，写像 $\symbf{x}\mapsto\symbf{A}\symbf{x}$ は $\mathbb{R}^n$ から $\mathbb{R}^m$ への線形写像である。$\symbf{A}$ の列ベクトルが一次独立であるとき，$\operatorname{rank}\symbf{A}=n$ である。正方行列 $\symbf{A}$ が正則であるための必要十分条件は $|\symbf{A}|\neq0$ であり，このとき $\symbf{A}\symbf{x}=\symbf{b}$ の解は $\symbf{x}=\symbf{A}^{-1}\symbf{b}$ と一意に定まる。$\symbf{E}$ は単位行列，${}^t\!\symbf{A}$ は転置行列を表す。
\end{bkhonbun}
\bkyaku{太字の $\symbf{A}$・$\symbf{x}$ は、この本の細字の行列 $A$ と矢印のベクトル $\vec x$。$\mathbb{R}^n$ は「$n$ 個の実数の組の全体」。「写像」「線形写像」は 7-2 の「入れると出てくる機械」で、線形はまっすぐはまっすぐのまま。「列ベクトルが一次独立」は、どの列もほかの列の組み合わせで作れないこと（7-4：2本の列が平行でない）。「rank」（階数）は、本当に独立な式の数（7-5）。「正則」は逆行列があること、$|\symbf{A}|$ は行列式 $\det A$（7-4）、「一意に定まる」は「ただ1つに決まる」。$\symbf{E}$ はこの本の $I$、${}^t\!\symbf{A}$ はこの本の $A^{\mathsf T}$（7-1。$A'$ と書く本もある）。}
\end{senmonsho}

\begin{bkcenter}
\begin{bktable}{colspec={Q[l,m,wd=34mm] Q[c,m,wd=36mm] Q[c,m,wd=42mm]}, rowsep=1pt, row{1}={font=\sffamily\small}, cells={font=\small}}
もの & この本 & 専門書でよく見る書き方 \\
行列 & $A$ & $\symbf{A}$、$[a_{ij}]$ \\
ベクトル & $\vec x$（たて書き） & $\symbf{x}$、$\symbfit{x}$ \\
単位行列 & $I$ & $E$、$\symbf{I}$、$\symbf{1}$ \\
転置 & $A^{\mathsf T}$ & ${}^t\!A$、$A'$、$A^\top$ \\
行列式 & $\det A$ & $|A|$ \\
固有値 & $\lambda$ & $\lambda$（制御の本では $s$ も） \\
四端子定数の行列 & 数で書いた2×2 & $\begin{pmatrix}A&B\\C&D\end{pmatrix}$、F 行列 \\
\end{bktable}
\bktabcaption{この本の記号と、専門書でよく見る記号の対応。専門書を読むときに見返そう。}
\end{bkcenter}


\kviilhead{この巻の道具と、働く場面}
第7巻の3つの部の道具が、この章のどの場面で働いていたかを、表にまとめます。1つの場面で、いくつもの道具がいっしょに働いていることがわかります。
\begin{bkcenter}
\begin{bktable}{colspec={Q[l,m,wd=40mm] Q[c,m,wd=26mm] Q[c,m,wd=26mm] Q[c,m,wd=28mm]}, rowsep=1pt, row{1}={font=\sffamily\footnotesize}, cells={font=\small}}
場面 & 行列の計算（7-1〜7-5） & 回転と座標（7-6・7-7） & 固有値とくり返し（7-8〜7-10） \\
ロボット（①②） & ● & ● & ● \\
ドローン・VR（③） & ● & ● & \\
回路・送電線（④〜⑥） & ● & & ● \\
建物・橋（⑦） & ● & ● & ● \\
画像・ゲーム（⑧） & ● & ● & ● \\
検索（⑨） & ● & & ● \\
GPS・データ（⑩⑪） & ● & & ● \\
制御（⑫） & ● & & ● \\
\end{bktable}
\bktabcaption{●は、その場面で使っている道具。「行列の計算」は、すべての場面の土台になっている。}
\end{bkcenter}

\begin{kotoba}
\small
\bkword{順運動学・逆運動学}{関節の角から手の位置を求めることと、手の位置から関節の角を求めること。}
\bkword{分布定数回路}{抵抗やコイル・コンデンサーのはたらきが、線全体に広がっている回路。長い電線。}
\bkword{特性インピーダンス}{長い線の性質を表す数 $Z_0=\sqrt{z/y}$。アンテナ線の「75\,Ω」など。}
\bkword{ページランク}{リンクの移り方の表の、固有値1の固有ベクトルでページの順位を決める方法。}
\end{kotoba}

\begin{furikaeri}
\begin{checklist}
\item ロボットの腕を、ヤコビ行列の逆行列で目標へ近づけるしくみを説明できる。\bkfurref{bk:7-11:gyaku}
\item ばねの網と抵抗の網、箱をつないだ回路を、行列で書ける。\bkfurref{bk:7-11:ami}
\item 長い電線が $\cosh$・$\sinh$ の行列になるわけを、箱のくさりとして読める。\bkfurref{bk:7-11:bunpu}
\end{checklist}
\end{furikaeri}

\bkfukidashi{sora}{表のまま計算する道具が、ロボットにも送電線にも検索にも入っていた。次の第8巻では、空間のあちこちに矢印がある「ベクトル場」を読みにいこう。}
